\section{Catálogo de características efectivas}\label{sec:catalogo}
Las fórmulas de la tabla \ref{tab:caracteristicas} son transcripción del contrato estructurado, antes de imputación y estandarización. $q_1,q_2,q_3$ denotan unidades de los tres meses anteriores. Las unidades monetarias son las de la extracción: no se infiere una moneda empresarial sin metadatos adicionales.
\begingroup\small
\begin{longtable}{r>{\raggedright\arraybackslash}p{6.7cm}>{\raggedright\arraybackslash}p{6.2cm}}\caption{Entradas efectivas y fórmulas del contrato. $q_j=Q_{p,t-j}$.}\label{tab:caracteristicas}\\\toprule Índice & Nombre & Fórmula registrada\\\midrule\endfirsthead\toprule Índice & Nombre & Fórmula registrada\\\midrule\endhead
1 & \texttt{\detokenize{unidades_lag_1}} & Q(p,t-1) \\
2 & \texttt{\detokenize{unidades_lag_2}} & Q(p,t-2) \\
3 & \texttt{\detokenize{unidades_lag_3}} & Q(p,t-3) \\
4 & \texttt{\detokenize{unidades_media_3}} & (q1+q2+q3)/3 \\
5 & \texttt{\detokenize{unidades_mediana_3}} & mediana(q1,q2,q3) \\
6 & \texttt{\detokenize{unidades_std_3}} & std(q1,q2,q3), ddof=0 \\
7 & \texttt{\detokenize{cambio_unidades_1}} & q1-q2 \\
8 & \texttt{\detokenize{variacion_relativa_1}} & (q1-q2)/q2 si q2>0; nulo en otro caso \\
9 & \texttt{\detokenize{tendencia_unidades_3}} & (q1-q3)/2 \\
10 & \texttt{\detokenize{facturas_lag_1}} & COUNT DISTINCT InvoiceID(p,t-1) \\
11 & \texttt{\detokenize{facturas_media_3}} & media de facturas distintas en t-1,t-2,t-3 \\
12 & \texttt{\detokenize{ingresos_lag_1}} & SUM VentaSinImpuesto(p,t-1) \\
13 & \texttt{\detokenize{ingresos_media_3}} & media de ingresos de t-1,t-2,t-3 \\
14 & \texttt{\detokenize{precio_unitario_ponderado_lag_1}} & ingresos\_lag\_1/q1 si q1>0; nulo si q1=0 \\
15 & \texttt{\detokenize{clientes_distintos_lag_1}} & COUNT DISTINCT NombreCliente(p,t-1) \\
16 & \texttt{\detokenize{unidades_por_factura_lag_1}} & q1/facturas\_lag\_1 si facturas\_lag\_1>0; nulo si no \\
17 & \texttt{\detokenize{antiguedad_observada_meses}} & ordinal(t)-ordinal(primer mes observado antes del corte) \\
18 & \texttt{\detokenize{mes_objetivo_sin}} & sin(2*pi*(mes(t)-1)/12) \\
19 & \texttt{\detokenize{mes_objetivo_cos}} & cos(2*pi*(mes(t)-1)/12) \\
20 & \texttt{\detokenize{dias_mes_objetivo}} & días calendario de t, con bisiestos \\
\bottomrule\end{longtable}
\endgroup

\section{Trazabilidad y reproducción documental}
\subsection{Relación entre afirmaciones y evidencia}
El archivo \archivo{informe/evidencias/trazabilidad.json} relaciona cada tabla y figura con su fuente y hash SHA-256. Las afirmaciones cuantitativas sobre convergencia se derivan de los CSV por época y del registro de inicialización; las de evaluación se contrastan con las matrices, predicciones y análisis estadístico. Los pesos se comparan exactamente con el estado del archivo de congelación. Las figuras copiadas mantienen los mismos bytes de su origen; las nuevas vistas de soporte temporal, comparación global con cinco métricas, convergencia, paneles de error y error por partición se generan desde registros guardados.

La procedencia principal se organiza así:
\begin{itemize}
\item Calidad, grano y cardinalidades: \archivo{resultados/auditoria/}, especialmente \archivo{perfilado_datos.json} y \archivo{comparacion_csv.json}.
\item Fórmulas, parámetros y particiones: \archivo{datos/Procesados/fase_03/} y su manifiesto.
\item Pesos y recorrido del perceptrón: \archivo{resultados/fase_04/modelo_seleccionado.json} e \archivo{historial_orden_barajado.csv}.
\item Logística, gradientes y selección: \archivo{resultados/fase_05/}, sus JSON y \archivo{historial_eta_010.csv}.
\item PCA y estados congelados: \archivo{resultados/fase_06/}.
\item Métricas y predicciones finales alineadas: \archivo{resultados/fase_07/}.
\item Descomposición temporal, acuerdos, distribuciones y sensibilidad agrupada: \archivo{resultados/fase_08/}.
\end{itemize}
Los nombres técnicos de directorio se mantienen para reproducibilidad; no organizan la narrativa académica.

\subsection{Comandos sin reentrenamiento}
Desde la raíz del proyecto, el exportador documental lee artefactos y escribe únicamente dentro de \archivo{informe/}:
\begin{verbatim}
python informe/herramientas/generar_informe.py
cd informe
pdflatex -interaction=nonstopmode -halt-on-error main.tex
pdflatex -interaction=nonstopmode -halt-on-error main.tex
pdflatex -interaction=nonstopmode -halt-on-error main.tex
\end{verbatim}
No ejecutar los entrenadores, la preparación del dataset ni la evaluación final para compilar. Los archivos \archivo{README.md} y \archivo{evidencias/verificacion_informe.json} documentan los comandos efectivos, número de tablas/figuras y controles de integridad. Se comprueban referencias cruzadas, ecuaciones, archivos gráficos y dimensiones de tablas, además de la inspección visual del PDF.

Los antecedentes registran 104 pruebas satisfactorias, sin fallos ni omisiones, antes de esta elaboración documental. El informe verifica nuevamente integridad y consistencia de las fuentes sin modificar esas pruebas o artefactos. Las pruebas matemáticas históricas usan ejemplos de control y no deben confundirse con reentrenamientos de WWI.

\subsection{Parámetros y series que no deben confundirse}
Los JSON conservan precisión completa de pesos y sesgos; los porcentajes del texto se redondean solo para lectura. Las curvas de accuracy y BA registran estados por época; las métricas finales corresponden al checkpoint restaurado. La norma del gradiente logístico incluye el sesgo y se registra después de cada actualización; \archivo{gradient_used_norm} corresponde al estado previo. No se transforma una curva de entrenamiento en una curva de test.

\section{Continuidad obligatoria del proyecto académico}
La siguiente ampliación documental deberá agregar, para cada Data Mart verificado, su formulación, procedencia, metodología, resultados y conclusiones parciales. No se presupone que Cliente o Producto constituyan Data Marts distintos; deben identificarse los procesos reales del DW y las extracciones disponibles. No se incluyen capítulos de resultados vacíos o ficticios. Antes de una entrega definitiva quedan pendientes la identificación del ciclo actual y la evidencia operacional de cobertura, claves y unidades. Esta versión termina con Ventas documentado y evaluado, manteniendo explícita la obligación de completar las demás aplicaciones.
